\providecommand{\kBKM}{\kappa_{\mathrm{BKM}}}
\providecommand{\kch}{\kappa_{\mathrm{ch}}}
\providecommand{\kcat}{\kappa_{\mathrm{cat}}}
\providecommand{\kfib}{\kappa_{\mathrm{fiber}}}
\providecommand{\CoHA}{\mathrm{CoHA}}
\providecommand{\Perf}{\mathrm{Perf}}
\providecommand{\Rep}{\mathrm{Rep}}
\providecommand{\SpCh}{\mathrm{Sp}^{\mathrm{ch}}}
\providecommand{\PhiFA}{\Phi^{\mathrm{FA}}}
\providecommand{\cT}{\mathcal{T}}
\providecommand{\cW}{\mathcal{W}}
\providecommand{\cO}{\mathcal{O}}
\providecommand{\CC}{\mathbb{C}}
\providecommand{\ZZ}{\mathbb{Z}}

\section{Class-$\mathcal S$ parent and genus-two obstruction}
\label{sec:k3e-class-s-parent}

The $K3\times E$ chiral-Borcherds bridge has three independent inputs.
The class-$\mathcal S$ anomaly source is the $A_1$ theory on the
twenty-four-punctured sphere $\Sigma_{0,24}$.  The genus-two coordinate
is the Siegel parameter of the Borcherds denominator.  The CY$_3$
functorial output is $E_1$-chiral; $E_2$ structure is carried by center,
double, and module-category constructions.

\subsection{Theorem: surface, curve, and Siegel roles}
\label{subsec:k3e-three-role-separation}

For $X=K3\times E$, the surface used in the Stage-$2$ specialization,
the class-$\mathcal S$ curve, and the genus-two automorphic coordinate
live in different categories:
\[
  \SpCh_{\Sigma_2,C}\bigl(\PhiFA_3(\Perf(X))\bigr),\qquad
  \cT[A_1,\Sigma_{0,24}],\qquad
  \Omega\in\mathbb H_2.
\]
Here $\Sigma_2\subset X$ is a surface input for the CY$_3$
specialization and $C\subset X$ is the reference curve on which the
$E_1$-chiral algebra is read.  The punctured sphere $\Sigma_{0,24}$
belongs to the four-dimensional class-$\mathcal S$ parent and supplies
the $SU(2)^{24}$ flavor datum.  The coordinate
$\Omega=(\tau,z,\sigma)$ belongs to the Siegel half-space on which the
Borcherds product $\Delta_5(\Omega)$ lives.  The bridge compares these
roles through protected indices and Borcherds lifts while keeping the
three geometric inputs distinct.

\subsection{Theorem: pants arithmetic on $\Sigma_{0,24}$}
\label{subsec:sigma-0-24-pants-arithmetic}

For an $A_1$ class-$\mathcal S$ theory on a genus-$g$ curve with $n$
maximal regular punctures, a pants decomposition has
\[
  N_T = 2g-2+n,\qquad N_G = 3g-3+n.
\]
Each $T_2$ trinion contributes $(n_v,n_h)=(0,4)$ and each gauged
$SU(2)$ tube contributes $(n_v,n_h)=(3,0)$ in the
Shapere-Tachikawa normalization.  At $(g,n)=(0,24)$,
\[
  N_T=22,\qquad N_G=21,
\]
hence
\[
  (n_v,n_h)_{\Sigma_{0,24}}
  =22(0,4)+21(3,0)
  =(63,88).
\]
The four-dimensional and Schur-sector two-dimensional central charges
are
\[
  c_{4d}\bigl(\cT[A_1,\Sigma_{0,24}]\bigr)
  =\frac{2n_v+n_h}{12}
  =\frac{126+88}{12}
  =\frac{107}{6},
\]
\[
  c_{2d}\bigl(\cT[A_1,\Sigma_{0,24}]\bigr)
  =-12\,c_{4d}
  =-214.
\]
The twenty-four maximal punctures give the $SU(2)^{24}$ flavor datum
for the K3 elliptic-genus and Mathieu-covariant
averaging comparison.

\subsection{Lemma: closed genus-two candidate fails}
\label{subsec:closed-genus-two-candidate-fails}

The closed genus-two $A_1$ curve has two $T_2$ trinions and three
gauged $SU(2)$ tubes:
\[
  N_T=2,\qquad N_G=3.
\]
The raw comparison tuple is
\[
  (N_G,n_h)_{\mathrm{closed}\ g=2}=(3,8),
\]
where $N_G$ counts gauge tubes.  The central-charge formula uses the
effective vector count $n_v^{\mathrm{eff}}=3N_G=9$, so
\[
  c_{4d}^{\mathrm{closed}\ g=2}
  =\frac{2n_v^{\mathrm{eff}}+n_h}{12}
  =\frac{18+8}{12}
  =\frac{13}{6},
\qquad
  c_{2d}^{\mathrm{closed}\ g=2}=-26.
\]
A closed curve has no puncture flavor group and carries zero
$SU(2)^{24}$ boundary factors.  Its anomaly pair $(13/6,-26)$ misses
the required pair $(107/6,-214)$, and its flavor datum misses the
twenty-four maximal punctures of $\cT[A_1,\Sigma_{0,24}]$.  The
closed genus-two candidate therefore fails both gates of the
$K3\times E$ bridge.

\subsection{Theorem: CY$_3$ output scope}
\label{subsec:k3e-cy3-output-scope}

Let $X=K3\times E$ and let
\[
  A_X=\Phi_3(\Perf(X)).
\]
Since $\dim_{\CC}X=3$, the specialized output $A_X$ is an $E_1$-chiral
object.  The $E_2$ layer is obtained from $Z(\Rep(A_X))$, the
Hall--Drinfeld double of the positive half, and the module-center
formalism.  A surface-level $E_2$ statement belongs to one of these
derived layers.  The Hall--Drinfeld/Borcherds branch uses the compact Hall
primitive and automorphic denominator.  The Mukai self-mirror branch
uses stable envelopes and the rank-\(24\) Mukai current algebra.

\subsection{Theorem: positive-half CoHA normalization}
\label{subsec:positive-half-coha-normalization}

The toric local model fixes the algebra-half convention:
\[
  \CoHA(\CC^3)=Y^+(\widehat{\mathfrak{gl}}_1).
\]
The vertex algebra $\cW_{1+\infty}$ appears after adjoining the
opposite half and Cartan modes, completing the Hall--Drinfeld double,
and evaluating along the chiral spectral line.  This is a construction
from an $E_1$ positive half to a vertex-algebraic object.

\subsection{Theorem: separated invariants of $K3\times E$}
\label{subsec:k3e-separated-invariants}

\[
\begin{array}{c|c|l}
\text{invariant} & \text{value} & \text{calculation}\\
\hline
\kcat(K3\times E) & 0 &
  \chi(\cO_{K3})\,\chi(\cO_E)=2\cdot 0\\
\kch(K3\times E) & 0 &
  \sum_q(-1)^qh^{0,q}=1-1+1-1\\
\kch^{\mathrm{Heis}}(K3\times E) & 3 &
  \text{Heisenberg-Mukai specialization}\\
\kBKM(\mathfrak g_{\Delta_5}) & 5 &
  \mathrm{wt}(\Delta_5)=c_1(0)/2=10/2\\
\kfib(K3\times E) & 24 &
  \operatorname{rk} H^\bullet(K3,\ZZ)
\end{array}
\]
These numbers are attached to distinct constructions.  The bridge uses
them as independent invariants, with the compact total-space values
separated from the Heisenberg and Borcherds witnesses.

\subsection{Comparison: Schur sector to Borcherds denominator}
\label{subsec:schur-borcherds-comparison}

The comparison has the form
\[
  \operatorname{Schur}\bigl(\cT[A_1,\Sigma_{0,24}]\bigr)
  \longrightarrow
  \phi^{K3}_{0,1}
  \xrightarrow{\operatorname{Borch}}
  \Delta_5.
\]
The first arrow is the Mathieu-covariant projection from the protected
sector to the K3 weak Jacobi form.  The second arrow is the
Borcherds lift with weight
\[
  \kBKM(\mathfrak g_{\Delta_5})=c_1(0)/2=5.
\]
Identifying the Schur VOA with the $\Delta_5$ chamber requires an
OPE-level construction beyond the central-charge arithmetic above:
the protected Schur algebra must be mapped to the completed
Hall--Drinfeld/Borcherds object with the positive half, negative half,
Cartan, pairing, and center data preserved.

\paragraph{References.}
Gaiotto 2012 for class-$\mathcal S$; Chacaltana-Distler 2010 for the
$A_1$ pants decomposition; Shapere-Tachikawa 2008 for the
four-dimensional central-charge normalization;
Beem-Lemos-Liendo-Peelaers-Rastelli-van Rees 2015 for
$c_{2d}=-12c_{4d}$; the Gritsenko-Nikulin 1998 Borcherds product
producing $\Delta_5$.
